\section{Materials and Methods}

\subsection{Study Area and Datasets}
The study covers sugarbeet cultivation districts in Germany, utilizing a multi-source dataset spanning 1981--2024. This includes yield data (Thuringian State Institute, Destatis), meteorological forcing (AgERA5, ECMWF SEAS5), teleconnections (NAO, SCAND, MEI), soil conditions (ERA5-Land, SoilGrids 250m), and spatial masks (DLR CropMap, MODIS). Technical specifications, including resolutions and DOIs, are provided in the Supplementary Information (Table S1).

\subsection{The Statistical Trend Baseline (GAM + ARIMA)}
The Statistical Trend Baseline model ($M_1$) is a robust time-series forecasting pipeline executed via strict Walk-Forward Validation. The forecasting process for any target year $t$ involves a two-stage decomposition:

\begin{enumerate}
    \item \textbf{Trend Component (Long-Term):} The secular technological trend is modeled using a Generalized Additive Model (GAM) where $f(t)$ is a smoothing spline (10 splines):
    \[ \hat{y}_{\text{trend}}(t) = f(t) \]

    \item \textbf{Residual Component (Short-Term):} Let $R_i = y_i - \hat{y}_{\text{trend}}(i)$ represent the historical residuals for years $i < t$. This sequence is modeled using an ARIMA(1, 0, 0) process to capture short-term autocorrelation:
    \[ \hat{y}_{\text{residual}}(t) = \text{ARIMA}(R_{1 \dots t-1}) \]
\end{enumerate}

The final trend forecast $M_1$ for year $t$ is the sum of these two components:
\[ M_1(t) = \hat{y}_{\text{trend}}(t) + \hat{y}_{\text{residual}}(t) \]

Following the recommendation of \citet{paudel2022field}, we employed a time-based expanding window for training and validation to avoid temporal leakage.

\subsection{Feature Engineering}
Our two-stage pipeline translates raw data into agronomic signals. First, a \textbf{Bio-Physical Risk Scanner} within WOFOST \cite{boogaard1998wofost} calculates specific risks like Oxygen Stress (Anoxia) and Harvest Respiration. Second, we compute rigorous Z-scores ($Z_t$) using an expanding window to normalize extreme events relative only to known history. Key indices include $Z_{\text{heat}}$ (heat stress), $Z_{\text{bal}}$ (water balance), $Z_{\text{tank}}$ (winter recharge), $Z_{\text{anoxia}}$ (oxygen stress), and $Z_{\text{sow}}$ (sowing delay), which are aggregated into compound indices (Scorch, Drown, LateStart) detailed in the Supplementary Information.

\subsection{Predictive Algorithms}
Machine learning components use Scikit-Learn \cite{scikitlearn} and XGBoost \cite{chen2016xgboost}. The system utilizes five specialized sub-models $M_k$:
\begin{itemize}
    \item \textbf{$M_1$: Statistical Trend Baseline} as defined in Section 2.2.
    \item \textbf{$M_2$: Native Ensemble}, a consensus-driven approach using weather and soil features via XGBoost and Random Forest.
    \item \textbf{$M_3$: Hybrid XGBoost}, which predicts the ratio between actual yield and WOFOST Stage 1 simulation, leveraging physics-based stress indices to correct simulation biases.
    \item \textbf{$M_4$: Robust Linear Integrator}, a HuberRegressor ($\epsilon=1.35$) predicting residuals from the statistical trend, providing robustness against non-climatic outliers.
    \item \textbf{$M_5$: V31 Solar-Gated Specialist}, utilizing conditional logic based on radiation flux and water balance to identify high-yield potential years.
\end{itemize}

\subsection{Super Ensemble Architecture}

The system integrates the five sub-models through an XGBoost-based Meta-Learner. The architecture follows a two-stage training and inference protocol:

\subsubsection{Cross-Validation Protocol}
Model validation employs a strict \textbf{expanding window (walk-forward)} strategy. For each year $t \in \{2000, \dots, 2024\}$, the Meta-Learner is trained exclusively on data from years $\{1981, \dots, t-1\}$.

\subsubsection{Meta-Learner Training (Discrete Selection)}
The Meta-Learner is trained as a multi-class classifier to predict the optimal sub-model $M_k$. For each training sample, we identify the sub-model that achieved the minimum absolute error, assigning it as the target label. To prioritize the mitigation of extreme errors, samples are weighted by the $\text{Regret}$ ($w_{d,t}$), defined as the difference between the median ensemble error ($E_{\text{median}}$) and the best-case error:
\begin{equation}
    w_{d,t} = \text{E}_{\text{median}}(d,t) - \min_{k \in \{1,\dots,5\}} |Y_{d,t} - \hat{y}_k(d,t)|
\end{equation}
While early iterations tested a 200~dt/ha outlier filter, the final architecture retains all data points, forcing the Meta-Learner to prioritize robustness through the weight $w_{d,t}$.

\subsubsection{Inference (Probabilistic Integration)}
During inference, the Meta-Learner generates a probability distribution $P(M_k | \mathbf{x}_{d,t})$ across all sub-models based on the feature vector $\mathbf{x}_{d,t}$ (comprising bio-physical Z-scores, sowing anomalies, and solar flux). The final yield forecast $\hat{Y}_{d,t}$ is calculated as the probabilistically weighted sum of all component predictions:
\begin{equation}
    \hat{Y}_{d,t} = \sum_{k=1}^5 P(M_k | \mathbf{x}_{d,t}) \cdot \hat{y}_k(d,t)
\end{equation}
This soft-voting approach allows the system to transition smoothly between specialized regimes, leveraging the Meta-Learner's confidence to blend mechanism-informed signals with statistical trends.

\begin{figure}[h]
    \centering
    \includegraphics[width=\linewidth]{super_ensemble_flowchart.png}
    \caption{Super Ensemble Architecture Flowchart. The system integrates multiple data sources and component models, using a Meta-Learner to dynamically integrate predictions based on the detected regime.}
    \label{fig:flowchart}
\end{figure}

\subsection{Uncertainty Quantification}
While the Super Ensemble produces point forecasts via probabilistic integration, component-level uncertainty is captured through quantile regression. Additionally, the entropy of the Meta-Learner's predicted probability distribution $P(M_k)$ provides an indirect measure of forecast confidence: low-entropy distributions signal high certainty in regime detection.